% concatenated from gnk2.tex
\documentclass{article}
\usepackage[utf8]{inputenc}
%\usepackage[russian]{babel}
\usepackage{amsmath,amsfonts,amsthm,amssymb}
\usepackage{graphicx}
\usepackage[all]{xy}


\newtheorem{lm}{Lemma}
\newtheorem{theorem}{Theorem}
\newtheorem{proposition}{Proposition}
\newtheorem{corollary}{Corollary}
\theoremstyle{definition}
\newtheorem{definition}{Definition}
\newtheorem{remark}{Remark}
\newtheorem{example}{Example}

\title{A generalisation of the Burau representation and groups $G_{n}^{3}$ for classical braids}

\author{Vassily Olegovich Manturov, Igor Mikhailovich Nikonov}
\date{}

\newcommand{\Z}{{\mathbb Z}}
\newcommand{\R}{{\mathbb R}}

\begin{document}

\maketitle

\abstract{We consider a certain modification of the group $G^3_n$ which describes dynamics of point configurations, in particular braids, and define a representation of the modified $G^3_n$. The braid representation induced is powerful enough to detect the kernel of the Burau representation.}

%\keywords
{Keywords: braid, $G^k_n$-group, Burau representation}

%\subclass
 {MSC2020: 57K10, 57K12, 20F36}





\section{Introduction}

In \cite{Ma} the first named author introduced a family of groups $G_{n}^{k}$ for two positive integers $n,k$ and formulated the following principle:

{\em If dynamical systems describing a motion of $n$ particles has a nice codimension $1$ property governed exactly by $k$
particles then these dynamical systems have a topological invariant valued in $G_{n}^{k}$.}

The first main example is an invariant of braids valued in $G_{n}^{3}$. 
Most of examples constructed over the last 10 years describe invariants of fundamental groups of configuration spaces
of $n$ points, which can be considered as ``braid-like objects'' that means an object with the number of particles fixed.

In this paper, we deal with a certain modification of the group $G^3_n$ which still describes dynamics of point configurations, in particular, braids. We define a representation of the modified group $\hat G^3_n$. The braid representation induced by it is powerful enough to detect the kernel of the Burau representation. 

\subsection{Acknowledgements}
The authors are grateful to Liu Yangzhou for valuable remarks. 

The first author was supported by the grant 25-21-00884 Applied combinatorial geometry and topology.



\section{A generalisation of the Burau representation}

Artin's presentation of the braid group $B_n$ has $n-1$ generators $\sigma_1,\dots,\sigma_{n-1}$ and the relations
\begin{itemize}
    \item (far commutativity) $\sigma_i\sigma_j=\sigma_j\sigma_i$ for any $|i-j|>1$;
    \item (Artin relations) $\sigma_i\sigma_{i+1}\sigma_i=\sigma_{i+1}\sigma_i\sigma_{i+1}$ for $i=1,\dots,n-2$.
\end{itemize}

Denote the natural homomorphism from $B_n$ to the permutation group $\Sigma_{n}$ on the set $\{1,\dots,n\}$ by $\rho$:
\begin{equation}\label{eq:map_from_braids_to_permutations}
\rho(\sigma_i)=(i\ i+1),\quad 1\le i\le n-1.
\end{equation}

The kernel $PB_n=\ker\rho$ is called the \emph{pure braid group} on $n$ strands.

Consider a modification of the groups $G^k_n$ introduced in~\cite[Section 2.3]{FKM}.

\begin{definition}
    The group $\hat G^3_n$ is the group with generators $a_{ijk}$ where $i,j,k$ are different numbers in $\{1,\dots,n\}$, and relations
\begin{enumerate}
    \item $a_{kji}=a_{ijk}^{-1}$;
    \item $a_{ijk}a_{pqr}=a_{pqr}a_{ijk}$ if there are at least 5 different numbers among $i,j,k,p,q,r$;
    \item $a_{ijk}a_{ijl}a_{ikl}a_{jkl}=a_{jkl}a_{ikl}a_{ijl}a_{ijk}$.
\end{enumerate}
\end{definition}

The permutation group $\Sigma_n$ acts on $\hat G^3_n$ by renumbering the indices:
\begin{equation}\label{eq:permutation_action_Gkn}
\tau \left(\prod_p a_{ijk}\right)=\prod_p a_{\tau(i)\tau(j)\tau(k)},\quad \tau\in \Sigma_n.
\end{equation}
This action defines a group structure on $\Sigma_n\ltimes \hat G^3_n$ by the formula
\begin{equation}\label{eq:product_Gkn_permutation}
(\sigma_1, \omega_1)\cdot(\sigma_2,\omega_2)=(\sigma_1\sigma_2, \sigma_2(\omega_1)\omega_2),\quad (\sigma_i,\omega_i)\in \Sigma_n\ltimes \hat G^3_n.
\end{equation}

\begin{figure}[h]
\centering\includegraphics[width=0.3\textwidth]{aijk.eps}
\caption{The event $a_{ijk}$}
\label{fig:aijk}
\end{figure}

In~\cite{MN}, a homomorphism from $PB_n$ to $G^3_n$ was constructed. To do that, one looks at the braids as loops in the configuration space of $n$-point subsets of the plane. More concretely, a braid $\beta$ is given by a continuous map $x(t)=(x_1(t),...,x_n(t))$, $t\in[0,1]$, where $\{x_1(t),...,x_n(t)\}$ is an $n$-point subset of $\R^2$ for any $t$, and $\{x_1(0),...,x_n(0)\}=\{x_1(1),...,x_n(1)\}\subset S^1$. The corresponding element $\phi_n(\beta)$ of $G^3_n$ is constructed as follows: let $0<t_1<\cdots<t_p<1$ be the moments when three points in the subset $x(t)$ become collinear and let $x_{i_r}(t_r), x_{j_r}(t_r), x_{k_r}(t_r)$ be those points for the moment $t_r$, $r=1,\dots,p$. Then
\[
\phi_n(\beta)=\prod_{r=1}^p a_{i_rj_rk_r}.
\]

The homomorphism $\phi_n\colon PB_n\to G^3_n$ can be generalised to a homomorphism from $B_n$ to $\Sigma_n\ltimes \hat G^3_n$.

\begin{proposition}
The map $\phi_n$ whose values in the generators $\sigma_i$, $i=1,\dots,n-1$, are
\begin{equation}\label{eq:map_from_braid_to_Gkn}
\phi_n(\sigma_i)=\left(\rho(\sigma_i), a_{i-1,i+1,i}a_{i-2,i+1,i}\cdots a_{1,i+1,i}\cdot a_{n,i+1,i}\cdots a_{i+2,i+1,i}\right).
\end{equation}
induces a well defined homomorphism $\phi_n\colon B_n\to\Sigma_n\ltimes \hat G^3_n$.
\end{proposition}

\begin{figure}[h]
\centering\includegraphics[width=0.3\textwidth]{sigmadynamics.pdf}
\caption{Dynamics of the generator $\sigma_i$}
\label{fig:sigmadynamics}
\end{figure}

\begin{proof}
By considering the dynamics for the Artin generators (Fig.~\ref{fig:sigmadynamics}), we get the formula~\eqref{eq:map_from_braid_to_Gkn}. The correctedness of the map $\phi_n$ is proved analogous to~\cite[Theorem 2]{Ma}.
\end{proof}

It can be said that the proof of invariance follows from the fact that THE GENERATORS CORRESPOND TO CODIMENSION 1, and THE RELATIONS CORRESPOND TO CODIMENSION 2.



\begin{corollary}
    The map  $\phi_n$ induces a homomorphism from $PB_n$ to $\hat G^3_n$.
\end{corollary}
\begin{proof}
 For any $\beta\in PB_n$, we have $\rho(\beta)=1$. Then $\phi_n$ maps $PB_n$ in the subgroup $1\times \hat G^3_n\subset \Sigma_n\ltimes \hat G^3_n$ which is isomorphic to $\hat G^3_n$.
\end{proof}

Let $A=\Z[t_1^\pm,\dots,t_n^\pm,s_1^\pm,\dots,s_n^\pm]$. Consider a free module 
\[
V=A\langle x_{ij}\mid i,j\in\{1,\dots,n\}, i\ne j\rangle
\]
of rank $n(n-1)$. Define an endomorphism $\rho_{ijk}\in End_A(V)$ using the formula
\[
\rho_{ijk}(x_{pq})=\left\{\begin{array}{cl}
    t_ix_{ij}+(1-t_i)x_{ik}, & (p,q)=(i,j), \\
    t_k^{-1}x_{kj}+(1-t_k^{-1})x_{ki}, & (p,q)=(k,j), \\
    s_j x_{jk}, & (p,q)=(j,k),\\
    s_j^{-1} x_{ji}, & (p,q)=(j,i),\\
    x_{pq}, & \mbox{otherwise}.
\end{array}\right.
\]

\begin{theorem}
    The map $a_{ijk}\mapsto \rho_{ijk}$ defines a representation of the group $\hat G^3_n$ on $V$.
\end{theorem}
\begin{proof}
    The relations 1 and 2 follow from the definition. The correctness of the relation 3 is verified directly (it is enough to consider the case $(i,j,k,l)=(1,2,3,4)$).
\end{proof}

\begin{corollary}
    The composition $\rho\circ\phi_n\colon PB_n\to End_A(V)$ is a representation of the pure braid group on the free module $V$.
\end{corollary}

\begin{example}
Let $n=5$ and $\beta\in PB_5$ be the element in the kernel of the Burau representation considered in~\cite{Big}. The braid $\beta$ is defined by the formula 
\[
\beta = [\psi_1^{-1}\sigma_4\psi_1,\psi_2^{-1}\sigma_4\sigma_3\sigma_2\sigma_1^2\sigma_2\sigma_3\sigma_4\psi_2],
\]
where $\psi_1=\sigma_3^{-1}\sigma_2\sigma_1^2\sigma_2\sigma_4^3\sigma_3\sigma_2$ and $\psi_2=\sigma_4^{-1}\sigma_3\sigma_2\sigma_1^{-2}\sigma_2\sigma_1^2\sigma_2^2\sigma_1\sigma_4^5$. Computer calculations show that the matrix $\rho\circ\psi_5(\beta)\in M_{20}(A)$ is not the identity. In particular, when $t_1=-1$ and $t_2=\cdots=t_5=s_1=\cdots=s_5=1$, the corner matrix element is
\[
\langle x_{12}|\rho\circ\psi_5(\beta)|x_{12}\rangle=-399.
\]
Thus, the constructed representation is not weaker than the Burau representation.
\end{example}

\begin{example}
    Consider the braid $\beta$ of the previous example as an element in the group $PB_6$. Liu Yangzhou~\cite{Yangzhou} proved that $\beta$ lies in the kernel of representation $\tilde\rho\circ\psi\circ f_d\circ p_6$, $d\ge 2$, considered in~\cite{MN23}. On the other hand, the matrix $\rho\circ\psi_6(\beta)\in M_{30}(A)$ is not the identity when $t_1=s_1=-1$ and $t_2=\cdots=t_6=s_2=\cdots=s_6=1$ (the corner matrix element is $\langle x_{12}|\rho\circ\psi_6(\beta)|x_{12}\rangle=-399$). Hence, the representation $\rho\circ\psi_6$ is stronger than the representation $\tilde\rho\circ\psi\circ f_d\circ p_6$.
\end{example}

\begin{remark}
    By specializing the values of the parameters $t_i,s_i$ in $\mathbb R$, we get a representation of the group $\hat G^3_n$ in the space $V_{\mathbb R}=\mathbb R\langle x_{ij}\mid i,j\in\{1,\dots,n\}, i\ne j\rangle$.  
    
    On the other hand, the space $V_{\mathbb R}$ splits into a direct sum of the subspaces $V_{sym}=\mathbb R\langle \frac{x_{ij}+x_{ji}}2\mid i\ne j\rangle$ and $V_{alt}=\mathbb R\langle \frac{x_{ij}-x_{ji}}2\mid i\ne j\rangle$. In~\cite[Section 2.3]{FKM}, several representations of $\hat G^3_n$ on $V_{sym}$ were defined. One can extend those representations to $V_{\mathbb R}$ by fixing a representation of $\hat G^3_n$ on $V_{alt}$ (for example, the trivial one). As can be seen, the representations obtained are different from those induced by the specializations of $\rho$. This indicates the existence of a more general representation that would encompass all of the above mentioned cases, but this representation has yet to be found. 
\end{remark}


\begin{thebibliography} {100}

\bibitem{Big} S. Bigelow, The Burau representation is not faithful for $n=5$, \emph{Geom. Topol.} {\bf 3} (1999) 397--404.

\bibitem{Ma}
V.O.Manturov, Non-Reidemeister knot theory and its applications in dynamical systems, geometry, and topology, arXiv:1501.05208.

\bibitem{MN}
V.O.Manturov, I.M.Nikonov, On braids and groups $G_{n}^{k}$, \emph{J. Knot Theory and Ramifications} {\bf 24}:13 (2015) 1541009.

\bibitem{Manturovbook} V.O. Manturov, \emph{Knot theory: Second edition}, CRC press, 2018.

\bibitem{FKMN} V.O. Manturov, D. Fedoseev, S. Kim, I. Nikonov,  {\em Invariants And Pictures: Low-dimensional Topology And Combinatorial Group Theory, Series On Knots And Everything} {\bf 66}, World Scientific, 2020.

\bibitem{FKM} D. Fedoseev, S. Kim, V.O. Manturov, Representations of $G^k_n$ and related groups, \emph{J. Knot Theory and Ramifications} {\bf 29}:13 (2020) 2043006.

\bibitem{FMN} D. Fedoseev, V.O. Manturov, I. Nikonov, Manifolds of triangulations, braid groups of manifolds, and the groups $\Gamma_n^k$, in: Garanzha, V.A., Kamenski, L., Si, H. (eds) \emph{Numerical Geometry, Grid Generation and Scientific Computing. Lecture Notes in Computational Science and Engineering} {\bf 143}, Springer, Cham, 2021, 13--36.

\bibitem{FKMN1} V.O. Manturov, D. Fedoseev, S. Kim, I. Nikonov, On groups $G_n^k$ and $\Gamma_n^k$: A study of manifolds, dynamics, and invariants, \emph{Bull. Math. Sci.} {\bf 11}:2 (2021) 2150004.

\bibitem{MN23}
V.O. Manturov, I.M. Nikonov, Maps to virtual braids and braid representations, \emph{Russian Math. Surveys}, {\bf 78}:2 (2023) 393--395; arXiv:2210.06862.

\bibitem{MN23}
V.O. Manturov, I.M. Nikonov, Maps from knots in the cylinder to flat-virtual knots, \emph{Russian Math. Surveys}, {\bf 79}:2 (2023) 366–368; arXiv:2210.09689.


\bibitem{Yangzhou} L. Yangzhou, private communication


%\bibitem{M22} V.O. Manturov, Realisability of $G^3_n$, realisability projection, and kernel of the $G^3_n$-braid presentation, arXiv:2210.13338.



\end{thebibliography}


\end{document}

\section{Introduction}

In \cite{Ma} the first named author introduced a family of groups $G_{n}^{k}$ for two positive integers $n,k$ and formulated the following principle:

{\em If dynamical systems describing a motion of $n$ particles has a nice codimension $1$ property governed exactly by $k$
particles then these dynamical systems have a topological invariant valued in $G_{n}^{k}$.}

The first main example is an invariant of braids valued in $G_{n}^{3}$. 
Most of examples constructed over the last 10 years describe invariants of fundamental groups of configuration spaces
of $n$ points, which can be considered as ``braid-like objects'' that means aan object with the number of particles fixed.

A very attractive question which we address in this paper is how to construct invariants of {\em knots} or {\em links},
objects which can not be described by a motion of fixed number of points, which do not form a group. 

To solve this question for links, we consider links as {\em equivalence classes of braids modulo Markov moves}.

The paper is organised as follows. In the next section we define the groups $G_{n}^{k}$, MN-indices and 
describe the main features how one easily gets invariants of braids valued in free product of some copies of the group
$\Z_{2}$. After that, we describe the construction of MN-indices for {\em multi-component links}, not only for braids.

Having this in mind, we represent links by equivalence classes of braids modulo Markov moves.

In section~\ref{sect:link_invariant}, we construct invariants of links.



\section{The groups $G_{n}^{k}$, braids, and $MN$-indices}

Artin's presentation of the braid group $B_n$ has $n-1$ generators $\sigma_1,\dots,\sigma_{n-1}$ and the relations
\begin{itemize}
    \item (far commutativity) $\sigma_i\sigma_j=\sigma_j\sigma_i$ for any $|i-j|>1$;
    \item (Artin relations) $\sigma_i\sigma_{i+1}\sigma_i=\sigma_{i+1}\sigma_i\sigma_{i+1}$ for $i=1,\dots,n-2$.
\end{itemize}

Denote the natural homomorphism from $B_n$ to the permutation group $\Sigma_{n}$ on the set $\{1,\dots,n\}$ by $\rho$:
\begin{equation}\label{eq:map_from_braids_to_permutations}
\rho(\sigma_i)=(i\ i+1),\quad 1\le i\le n-1.
\end{equation}

The kernel $PB_n=\ker\rho$ is called the \emph{pure braid group} on $n$ strands.


%\begin{definition}\label{def:gnk}
The groups of {\em free $k$-braids} $G_{n}^{k}$ ($n>k\geq 1$) are defined as groups with the set of generators $a_{m}$ which are indexed by all $k$-element subsets of $\{1,\ldots, n\}$,  and relations
 \begin{enumerate}
\item[(1)]
$(a_{m})^{2}=1$ for any unordered sets $m \subset \{1,\ldots, n\}, |m|=k$;
\item[(2)] (far commutativity) $a_{m}a_{m'} = a_{m'}a_{m},$ if $|m\cap m'| < k-1$;
\item[(3)] (tetrahedron relation) for every (ordered) subset $U=\{i_1,\dots,i_{k+1}\}$ in $\{1,\dots, n\}$ of cardinality $(k+1)$, let us denote them by $m^{p}=U\setminus\{i_p\}, p=1,\dots,k+1$. Then
$$a_{m^{1}}\cdots  a_{m^{k+1}} = a_{m^{k+1}}\cdots  a_{m^{1}}.$$
\end{enumerate}
%\end{definition}

The permutation group $\Sigma_n$ acts on $G^k_n$ by renumbering the indices:
\begin{equation}\label{eq:permutation_action_Gkn}
\sigma \left(\prod_p a_{i_{p1}i_{p2}\cdots i_{pk}}\right)=\prod_p a_{\sigma(i_{p1})\sigma(i_{p2})\cdots\sigma(i_{pk})}.
\end{equation}

In~\cite{MN}, a homomorphism from $PB_n$ to $G^3_n$ was constructed. To do that, one looks at the braids as loops in the configuration space of $n$-point subsets of the plane. More concretely, a braid $\beta$ is given by a continuous map $x(t)=(x_1(t),...,x_n(t))$, $t\in[0,1]$, where $\{x_1(t),...,x_n(t)\}$ is an $n$-point subset of $\R^2$ for any $t$, and $\{x_1(0),...,x_n(0)\}=\{x_1(1),...,x_n(1)\}$. The corresponding element $\phi_n(\beta)$ of $G^3_n$ is constructed as follows: let $0<t_1<\cdots<t_p<1$ be the moments when three points in the subset $x(t)$ become collinear and let $x_{i_r}(t_r), x_{j_r}(t_r), x_{k_r}(t_r)$ be those point for the moment $t_r$, $r=1,\dots,p$. Then
\[
\phi_n(\beta)=\prod_{r=1}^p a_{i_rj_rk_r}.
\]

The homomorphism $\phi_n\colon PB_n\to G^3_n$ defined in~\cite{MN} extends to a homomorphism $\phi_n\colon B_n\to\Sigma_n\ltimes G^3_n$ where the group structure on $\Sigma_n\ltimes G^3_n$ is defined by the formula
\begin{equation}\label{eq:product_Gkn_permutation}
(\sigma_1, \omega_1)\cdot(\sigma_2,\omega_2)=(\sigma_1\sigma_2, \sigma_2(\omega_1)\omega_2).
\end{equation}

The values of the homomorphism $\phi_n$ on the generators $\sigma_i$, $i=1,\dots,n-1$, are %(Fig.~\ref{fig:sigmadynamics})
\begin{equation}\label{eq:map_from_braid_to_Gkn}
\phi_n(\sigma_i)=\left(\rho(\sigma_i), a_{i-1,i,i+1}a_{i-2,i,i+1}\cdots a_{1,i,i+1}\cdot a_{i,i+1,n}a_{i,i+1,n-1}\cdots a_{i,i+1,i+2}\right).
\end{equation}

\begin{figure}[h]
\centering\includegraphics[width=0.3\textwidth]{sigmadynamics.pdf}
\caption{Dynamics of the generator $\sigma_i$}
\label{fig:sigmadynamics}
\end{figure}

Let us recall the definition of MN-indices.

Given a generator $a_m$ of $G^k_n$ and an index $l\not\in m$, consider the group 
\[
Z_{m;l}=\bigoplus_{i\in m}\Z_2 e_i/(\sum_{i\in m}e_i=0)\simeq \Z_2^{\oplus k-1}
\]
and the homomorphism $\psi_{m;l}\colon G^k_n\to Z_{m;l}$ such that $\psi_{m;l}(a_{m\setminus i\cup l})=e_{i}$, and $\psi_{m;l}(a_{m'})=0$ for the other generators of $G^k_n$.

Denote $Z_{m}=\bigoplus_{l\not\in m}Z_{m;l}$ and $\psi_{m}=\bigoplus_{l\not\in m}\psi_{m;l}\colon G^k_n\to Z_{m}$. 

Let $H_{m}=\Z_2^{\ast Z_{m}}$ be the free product of the groups $\Z_2$ labeled by the elements of $Z_{m}$. Denote the generators of $H_{m}$ by $f_z, z\in Z_{m}$. 

The group $Z_{m}$ acts on $H_{m}$ by the formula
\begin{equation}\label{eq:Z_action_H}
z\cdot\left(\prod_p f_{z_p}\right)=\prod_p f_{z_p+z}, \quad z,z_p\in Z_{m}.
\end{equation}

There is a homomorphism $\tilde\psi_{m}\colon G^k_n\to Z_{m}\ltimes H_{m}$ where the multiplication of $Z_{m}\ltimes H_{m}$ is defined by the formula
\begin{equation}\label{eq:ZH_product}
(z_1,h_1)\cdot (z_2,h_2) = (z_1+z_2, (z_2\cdot h_1)h_2), \quad z_1,z_1\in Z_{m}, h_1,h_2\in H_{m}.
\end{equation}
The values of the homomorphism $\tilde\psi_{m}$ on the generators $a_{m'}$ are  
\begin{equation}\label{eq:map_Gkn_to_ZH}
\tilde\psi_{m}(a_{m'})=\left\{\begin{array}{cc}
   (0, f_0),  & m'=m, \\
    (\psi_{m}(a_{m'}),1), & \mbox{otherwise}.  
\end{array}\right.
\end{equation}

%Given a generator $a_{ijk}$ of $G^3_n$ and an index $l\not\in\{i,j,k\}$, consider the group 
%\[
%Z_{ijk;l}=\Z_2\langle e_{ijl}, e_{ikl}, e_{jkl}\rangle/(e_{ijl}+e_{ikl}+e_{jkl}=0)\simeq \Z_2\oplus\Z_2
%\]
%and the homomorphism $\psi_{ijk;l}\colon G^3_n\to Z_{ijk;l}$ such that $\psi_{ijk;l}(a_{jkl})=e_{jkl}$, $\psi_{ijk;l}(a_{ikl})=e_{ikl}$, $\psi_{ijk;l}(a_{ijl})=e_{ijl}$, and $\psi_{ijk;l}(a)=0$ for the other generators of $G^3_n$.

%Denote $Z_{ijk}=\bigoplus_{l\not\in\{i,j,k\}}Z_{ijk;l}$ and $\psi_{ijk}=\bigoplus_{l\not\in\{i,j,k\}}\psi_{ijk;l}\colon G^3_n\to Z_{ijk}$. Let $H_{ijk}=\Z_2^{\ast Z_{ijk}}$ be the free product of the groups $\Z_2$ labeled by the elements of $Z_{ijk}$. Denote the generators of $H_{ijk}$ by $f_z, z\in Z_{ijk}$. The group $H_{ijk}$ acts on $H_{ijk}$ by the formula
%\[
%z\cdot\left(\prod_p f_{z_p}\right)=\prod_p f_{z_p+z}, \quad z,z_p\in Z_{ijk}.
%\]

%There is a homomorphism $\tilde\psi_{ijk}\colon G^3_n\to Z_{ijk}\ltimes H_{ijk}$ where the multiplication of $Z_{ijk}\ltimes H_{ijk}$ is defined by the formula
%\[
%(z_1,h_1)\cdot (z_2,h_2) = (z_1+z_2, (z_2\cdot h_1)h_2), \quad z_1,z_1\in Z_{ijk}, h_1,h_2\in H_{ijk}.
%\]
%The values of the homomorphism $\tilde\psi_{ijk}$ on the generators $a_{i'j'k'}$ are  
%\[
%\tilde\psi_{ijk}(a_{i'j'k'})=\left\{\begin{array}{cc}
%   (0, f_0),  & \{i',j',k'\}=\{i,j,k\}, \\
%    (\psi_{ijk}(a_{i'j'k'}),1), & \mbox{otherwise}.  
%\end{array}\right.
%\]

Denote $ZH=\prod_{m}Z_{m}\ltimes H_{m}$ and $\tilde\psi=\prod_{m}\tilde\psi_{m}\colon G^k_n\to ZH$. Any permutation $\sigma\in\Sigma_n$ induces homomorphisms $Z_{m;l}\to Z_{\sigma(m);\sigma(l)}$ by the formula
$\sigma(\sum_{i\in m}\lambda_{i}e_{i}) = \sum_{i\in m}\lambda_{i}e_{\sigma(i)}$,
homomorphisms $Z_{m}\to Z_{\sigma(m)}$ by the formula
\begin{equation}\label{eq:permutation_action_Z}
\sigma\left(\sum_{l\not\in m}z_{m;l}\right)=\sum_{l'\not\in\sigma(m)}\sigma(z_{m;\sigma^{-1}(l')}),
\end{equation}
and homomorphisms $H_{m}\to H_{\sigma(m)}$ by the formula
\[
\sigma\left(\prod_p f_{z_p}\right)=\prod_p f_{\sigma(z_p)}.
\]
These homomorphisms induce an action of the group $\Sigma_n$ on the group $ZH$:
\begin{equation}\label{eq:permutaion_action_ZH}
\sigma\cdot \prod_{m}\left(z_{m}, h_{m}\right)=\prod_{m}\left(\sigma(z_{\sigma^{-1}(m)}), \sigma(h_{\sigma^{-1}(m)})\right).
\end{equation}

Then the homomorphism $\tilde\psi$ extends to the homomorphism 
\[
\psi\colon \Sigma_n\ltimes G^k_n\to\Sigma_n\ltimes ZH
\]
defined by the formula
\[
\psi(\sigma,w)=(\sigma,\tilde\psi(w)).
\]

%Denote $ZH=\prod_{i,j,k}Z_{ijk}\ltimes H_{ijk}$ and $\tilde\psi=\prod_{i,j,k}\tilde\psi_{ijk}\colon G^3_n\to ZH$. Any permutation $\sigma\in\Sigma_n$ induces homomorphisms $Z_{ijk;l}\to Z_{\sigma(i)\sigma(j)\sigma(k);\sigma(l)}$ by the formula
%\[
%\sigma(\lambda_{1}e_{ijl}+\lambda_{2}e_{ikl}+\lambda_{3}e_{jkl}) = \lambda_{1}e_{\sigma(i)\sigma(j)\sigma(l)}+\lambda_{2}e_{\sigma(i)\sigma(k)\sigma(l)}+\lambda_{3}e_{\sigma(j)\sigma(k)\sigma(l)},
%\]
%homomorphisms $Z_{ijk}\to Z_{\sigma(i)\sigma(j)\sigma(k)}$ by the formula
%\[
%\sigma\left(\sum_{l\not\in\{i,j,k\}}z_{ijk;l}\right)=\sum_{l'\not\in\{\sigma(i),\sigma(j),\sigma(k)\}}\sigma(z_{ijk;\sigma^{-1}(l')}),
%\]
%and homomorphisms $H_{ijk}\to H_{\sigma(i)\sigma(j)\sigma(k)}$ by the formula
%\[
%\sigma\left(\prod_p f_{z_p}\right)=\prod_p f_{\sigma(z_p)}.
%\]
%These homomorphisms induce an action of the group $\Sigma_n$ on the group $ZH$:
%\[
%\sigma\cdot \prod_{i,j,k}\left(z_{ijk}, h_{ijk}\right)=\prod_{i,j,k}\left(\sigma(z_{\sigma^{-1}(i)\sigma^{-1}(j)\sigma^{-1}(k)}), \sigma(h_{\sigma^{-1}(i)\sigma^{-1}(j)\sigma^{-1}(k)})\right).
%\]

%Then the homomorphism $\tilde\psi$ extends to the homomorphism $\psi\colon \Sigma_n\ltimes G^3_n\to\Sigma_n\ltimes ZH$ defined by the formula
%\[
%\psi(\sigma,w)=(\sigma,\tilde\psi(w)).
%\]

\subsection{Component reduction of MN-indices}

Define a \emph{component map} as a map $c\colon\{1,\dots,n\}\to C$ from the set of indices to some set $C$ called the \emph{set of components}. Given a permutation $\sigma$, we say that a component map $c$ is \emph{compatible} with $\sigma$ if $c=c\circ\sigma$. A component map $c$ is compatible with a braid $\beta\in B_n$ if $c$ is compatible with the permutation $\rho(\beta)$.

For a braid $\beta$ one can define the \emph{canonical component map} 
\[
c_\beta=c_{\rho(\beta)}\colon\{1,\dots,n\}\to\{1,\dots,n\}/\rho(\beta)
\]
to the set of orbits (i.e. the cycles) of the permutation $\rho(\beta)$. Obviously, the component map $c_\beta$ is compatible with $\beta$.

Note that if $c$ is a components map compatible with a braid $\beta$ and $\beta'=\alpha\beta\alpha^{-1}$ is an adjoint braid then the component map $c'=c\circ\rho(\alpha)^{-1}$ is compatible with $\beta'$.

\begin{remark}
Consider the following general construction. Let $f\colon I\to A$ be a map, and $g\colon I\to J$, $h\colon A\to B$ other maps. 
%\[
%\xymatrix{I \ar[r]^f\ar[d]_g & A\ar[d]^h \\  J\ar@{.>}[r] & B}
%\]
%
The \emph{$(g,h)$-reduction} of the map $f$ is the map $f^{(g,h)}\colon J\to\Z[B]$ defined by the formula
\[
f^{(g,h)}(j)=\sum_{i\in g^{-1}(j)}h(f(i)).
\]

If $J$ is finite then there is a map 
\[
\prod_{i\in I} A\simeq Map(I,A)\to Map(J,\Z[B])\simeq \bigoplus_{j\in J}\Z[B],\quad f\mapsto f^{(g,h)}.
\]
\end{remark}

Now, apply the construction above to the $MN$-index map and the component map.

Let $C$ be a component set. For $s\subset C$, $|s|=k$, and $c\not\in s$, groups $Z^C_{s;c}$, $Z^C_s$, $H^C_s$ are defined analogously to $Z_{m;l}$, $Z_m$ and $H_m$. 

The component map $c$ induces homomorphisms $p^c_{m;l}\colon Z_{m;l}\to Z^C_{c(m);c(l)}$ when  $c|_{m\cup l}$ is injective (otherwise, we set $p^c_{m;l}=0$). Then one gets homomorphisms $p^c_m\colon Z_m\to Z^C_{c(m)}$, $p^c_{m}=\sum_{l\not\in m} p^c_{m;l}$, and $p^c_{m}\colon Z_{m}\ltimes H_{m}\to Z^C_{c(m)}\ltimes H^C_{c(m)}$.

Denote $A_C=\bigoplus_{s}\Z_2[Z^C_{s}\ltimes H^C_{s}]$.
The homomorphisms $p^c_{m}$ induce a map 
\[
p^c\colon ZH\to A_C
\]
where
\begin{equation}\label{eq:component_reduction}
p^c\left(\prod_{m}(z_{m},h_{m})\right)=\sum_{m}p^c_{m}(z_{m},h_{m}).
\end{equation}

Extend the map $p^c$ to a map $\Sigma_n\ltimes ZH\to A_C$ by the formula $p^c(\sigma,w)=p^c(w)$.

Let $V\subset A_C$ be the subspace generated by the elements
\[
p^{c\circ\sigma_1^{-1}}(y_1y_2y_1^{-1})-p^c(y_2),\quad y_i=(\sigma_i,w_i)\in \Sigma_n\ltimes ZH,\ i=1,2,
\]
and $\bar A_C=A_C/V$ be the quotient space. Then $p^c$ induces a map $\bar p^c\colon ZH\to \bar A_C$.

The composition 
\[
\bar p^c\circ\tilde\psi\colon G^k_n\to \bar A_C
\]
is called the \emph{component reduction of MN-indices map} induced by the component map $c$.


%Let $C$ be the component set of a component map. For distinct $u_1,u_2,u_3,u_4\in C$, consider the group 
%\[
%Z^C_{u_1u_2u_3;u_4}=\Z_2\langle e_{u_1u_2u_4}, e_{u_1u_3u_4}, e_{u_2u_3u_4}\rangle/(e_{u_1u_2u_4}+e_{u_1u_3u_4}+e_{u_2u_3u_4}=0).
%\]
%For distinct $u_1,u_2,u_3$, consider the groups $Z^C_{u_1u_2u_3}=\bigoplus_{u_4\not\in\{u_1,u_2,u_3\}}Z^C_{u_1u_2u_3;u_4}$, $H^C_{u_1,u_2,u_3}=\Z_2^{\ast Z^C_{u_1u_2u_3}}$ and $Z^C_{u_1u_2u_3}\ltimes H^C_{u_1,u_2,u_3}$, which are direct analogues of the groups $Z_{ijk}$, $H_{ijk}$ and $Z_{ijk}\ltimes H_{ijk}$.

%The component map $c\colon\{1,\dots,n\}\to C$ induces homomorphisms 
%\[
%p^c_{ijk;l}\colon Z_{ijk;l}\to Z^C_{c(i)c(j)c(k);c(l)},
%\]
%\[
%p^c_{ijk}\colon Z_{ijk}\to Z^C_{c(i)c(j)c(k)},\quad p^c_{ijk}=\sum_{l\not\in\{i,j,k\}}p^c_{ijk;l},
%\]
%\[
%p^c_{ijk}\colon H_{ijk}\to H^C_{c(i)c(j)c(k)}, \quad p^c_{ijk}\left(\prod_t f_{z_t}\right)=\prod_t f_{p^c_{ijk}(z_t)},
%\]
%and $p^c_{ijk}\colon Z_{ijk}\ltimes H_{ijk}\to Z^C_{c(i)c(j)c(k)}\ltimes H^C_{c(i)c(j)c(k)}$.

%Denote $A_C=\bigoplus_{c_1,c_2,c_3}\Z_2[Z^C_{c_1c_2c_3}\ltimes H^C_{c_1c_2c_3}]$.
%The homomorphisms $p^c_{ijk}$ induce a map 
%\[
%p^c\colon ZH\to A_C
%\]
%where
%\[
%p^c\left(\prod_{i,j,k}(z_{ijk},h_{ijk})\right)=\sum_{i,j,k}p^c_{ijk}(z_{ijk},h_{ijk}).
%\]

%Extend the map $p^c$ to a map $\Sigma_n\ltimes ZH\to A_C$ by the formula $p^c(\sigma,w)=p^c(w)$.

%Let $V\subset A_C$ be the subspace generated by the elements
%\[
%p^{c\circ\sigma_1^{-1}}(y_1y_2y_1^{-1})-p^c(y_2),\quad y_i=(\sigma_i,w_i)\in \Sigma_n\ltimes ZH,\ i=1,2,
%\]
%and $\bar A_C=A_C/V$ be the quotient space. Then $p^c$ induces a map $\bar p^c\colon ZH\to \bar A_C$.

%The composition 
%\[
%\bar p^c\circ\tilde\psi\colon G^3_n\to \bar A_C
%\]
%is called the \emph{component reduction of MN-indices map} induced by the component map $c$.

\section{MN-indices for arbitrary links in $3$-space}\label{sect:link_invariant}

Let $\beta\in B_n$ be a braid. Consider a geometric realization $x_i\colon [0,1]\to\R^2$, $i=1,\dots,n$, of the braid $\beta$ such that the points $x_1(0),\dots,x_n(0)$ lie in the unit circle of the plane. 

Extend the braid $\beta$ by adding $n$ strands $x_{i}\colon [0,1]\to\R^2$, $i=n+1,\dots,2n$, such that $x_i(t)\equiv x_i(0)$ for all $t\in[0,1]$ and $n+1\le i\le 2n$, and $x_{n+1}(0),\dots, x_{2n}(0)$ lie in the unit circle. We assume that the points $x_1(0),\dots, x_{2n}(0)$ lie in the cyclic order in the circle. Then there is a line $L\subset\R^2$ which separates the points $x_1(0),\dots,x_n(0)$ and the points $x_{n+11}(0),\dots,x_{2n}(0)$. We assume that the strands $x_i(t)$, $i\le n$, don't intersect the line $L$.

We call the obtained braid $\tilde\beta\in B_{2n}$ the \emph{closure} of the braid $\beta$.

Note that the union of the braid $\tilde\beta=\{x_i(t)\}_{1\le i\le n}$ and the segments $[x_i(0),x_{2n+1-i}(0)]$ and $[x_i(1),x_{2n+1-i}(1)]$, $i=1,\dots,n$, is the closure link of the braid $\beta$ (Fig.~\ref{fig:braid_closure}).

\begin{figure}[h]
\centering\includegraphics[width=0.8\textwidth]{braid_closure.pdf}
\caption{Closure of a braid}
\label{fig:braid_closure}
\end{figure}

Given a component map $c\colon\{1,\dots,n\}\to C$, consider its closure $\tilde c\colon\{1,\dots,2n\}\to C$ defined by the formula
\[
\tilde c(i)=\left\{\begin{array}{cc}
    c(i), & i\le n, \\
    c(2n+1-i), & i>n.
\end{array}\right.
\]

Define a map $w\colon B_n\to \bar A_C$ by the formula
\[
w(\beta)=\bar p^{\tilde c_{\beta}}\circ\tilde\psi\circ\phi(\tilde\beta).
\]

Note that if $x=\tilde\psi\circ\phi(\tilde\beta)$ then $x=\prod_{i<j<k}(z_{ijk},h_{ijk})$ where $z_{ijk}\in Z_{ijk}$ and $h_{ijk}=\prod_{t}f_{z_t}\in H_{ijk}$, $z_t\in Z_{ijk}$. In this case,
\[
w(\beta)=\bar p^{\tilde c_{\beta}}(x)=\prod_{c_1<c_2<c_3}\sum_{p(i)=c_1,p(j)=c_2,p(k)=c_3}(p^c(z_{ijk}),p^c(h_{ijk}))
\]
where $p^c(h_{ijk})=\prod_{t}f_{p^c(z_t)}$ and for any $e_\alpha\in Z_{ijk;p}$, $\alpha\in\{i,j,k\}$,
\[
p^c(e_\alpha)=\left\{\begin{array}{cl}
    e_{p(\alpha)}\in Z^C_{c_1,c_2,c_3; c(p)}, & c_1,c_2,c_3,c(p)\mbox{ are different} \\
     0, & \mbox{otherwise}.
\end{array}\right.
\]

\begin{thm}
    The map $w$ is invariant under first and second Markov moves.
\end{thm}

\begin{proof}
    1. Consider a first Markov move $\beta\mapsto\beta'=\beta_1\beta\beta_1^{-1}$. Denote $y=\psi\circ\phi(\tilde\beta)$, $y_1=\psi\circ\phi(\tilde\beta_1)\in\Sigma_n\ltimes ZH$. Then
\[
w(\beta')-w(\beta)=p^{c\circ\sigma_1^{-1}}(y_1yy_1^{-1})-p^c(y)=0\in\bar A_C,    
\]
where $c=\tilde c_\beta$ and $\sigma_1=\rho(\tilde\beta_1)$.

2. Let $\beta\mapsto \beta'=\sigma_n\beta\in B_{n+1}$ be a second Markov move. Then $c_{\beta'}(i)=c_\beta(i)$ for $i\le n$ and $c_{\beta'}(n+1)=c_\beta(n)$.

Let $v=\phi(\tilde\beta)\in G^3_{2n}$ and $v'=\phi(\tilde\beta')\in G^3_{2n+2}$. Then $v'=v'_1v'_2$ where $v'_1$ corresponds to $\sigma_n$ and $v'_2$ corresponds to $\beta$.

Since the $v'_1$ consists of generators $a_{n,n+1,j}$ and the impact of a generator $a_{ijk}$ to the image of the map $\psi$ is nontrivial only if the components $c(i)$, $c(j)$ and $c(k)$ are all different, we can ignore the term $v'_1$.

The word $v'_2$ differs from the word $v$ by the shift by $1$ of the indices which are greater than $n$, and by 2-letter subwords $a_{ij,n+1}a_{ij,n+2}$ which may appear in some places of $v'_2$. Since the strands $n+1$ and $n+2$. belong to one component, the impact of these generators in $Z_{c_1c_2c_3}$ by the map $p^c$ annihilates. Hence, $p^c(v)=p^c(v'_2)=p^c(v')$.
Thus, $w(\beta)=w(\beta')$.
\end{proof}

\begin{remark}
    In fact, $w$ is a link-homotopy invariant.
\end{remark}


\begin{example}
    Consider the braid $\beta=\sigma_3^2\in B_4$. Its closure link $L$ is the disjoint union of the Hopf link and two trivial components. Let us calculate the invariant value $w(L)$.

For the closure $\bar\beta$, we have 
\[
\phi(\bar\beta)=(a_{345}a_{346}a_{347}a_{348}a_{341}a_{342})^2.
\]
Then $\bar\psi_{345}(\phi(\bar\beta))=f_0f_{z_1+z_2+z_6+z_7+z_8}$ where $z_i=\psi_{345;i}(a_{34i})\in Z_{345;i}\subset Z_{345}$. Analogously, one computes $\bar\psi_{34k}(\phi(\bar\beta))$, $k=1,2,6,7,8$.

Let $\bar c$ be the closure of the canonical component map of the braid $\beta$. Then the component reduction $p^{\bar c}_{345}(\bar\psi_{345}(\phi(\bar\beta)))$ is zero because some of indices $3,4,5$ have same components: $\bar c(4)=\bar c(5)=4$. Analogously, $p^{\bar c}_{346}(\bar\psi_{346}(\phi(\bar\beta)))=0$.

On the other hand, 
\[
p^{\bar c}_{347}(\bar\psi_{347}(\phi(\bar\beta)))=p^{\bar c}_{347}(f_{z_5+z_6}f_{z_1+z_2+z_8})=f_0f_0=1\in Z^C_{234},
\]
because $\bar c(3)=3$, $\bar c(4)=4$, $\bar c(7)=2$ and $p^{\bar c}(z_5)=p^{\bar c}(z_6)=p^{\bar c}(z_2)=0$, $p^{\bar c}(z_1)=p^{\bar c}(z_8)$. Here $z_i$ denotes $\psi_{347;i}(a_{34i})\in Z_{347}$. Analogously,
$p^{\bar c}_{348}(\bar\psi_{348}(\phi(\bar\beta)))=p^{\bar c}_{134}(\bar\psi_{134}(\phi(\bar\beta)))=1\in Z^C_{134}$ and $p^{\bar c}_{234}(\bar\psi_{234}(\phi(\bar\beta)))=1\in Z^C_{234}$. Hence,
\begin{multline*}
p^{\bar c}(\bar\psi(\phi(\bar\beta)))=\\
p^{\bar c}_{134}(\bar\psi_{134}(\phi(\bar\beta)))+p^{\bar c}_{234}(\bar\psi_{234}(\phi(\bar\beta)))+p^{\bar c}_{347}(\bar\psi_{347}(\phi(\bar\beta)))+p^{\bar c}_{348}(\bar\psi_{348}(\phi(\bar\beta)))=\\
0\in A_C.
\end{multline*}
Thus, $w(\beta)=0$.
\end{example}



\begin{thebibliography} {100}

\bibitem{Ma}
V.O.Manturov, {\it
Non-Reidemeister knot theory and its applications in dynamical systems, geometry, and topology,} arXiv:1501.05208v1 [math.GT] 21 Jan 2015.

\bibitem{MN}
V.O.Manturov, I.M.Nikonov, On braids and groups $G_{n}^{k}$, \emph{J. Knot Theory and Ramifications} {\bf 24}:13 (2015) 1541009.

\bibitem{FKMN} V.O. Manturov, D. Fedoseev, S. Kim, I. Nikonov,  {\em Invariants And Pictures: Low-dimensional Topology And Combinatorial Group Theory, Series On Knots And Everything} {\bf 66}, World Scientific, 2020.

\bibitem{Manturovbook} V. O. Manturov Knot theory: Second edition, CRC press, 2018.

\bibitem{FMN} D. Fedoseev, V.O. Manturov, I. Nikonov, Manifolds of triangulations, braid groups of manifolds, and the groups $\Gamma_n^k$, in: Garanzha, V.A., Kamenski, L., Si, H. (eds) \emph{Numerical Geometry, Grid Generation and Scientific Computing. Lecture Notes in Computational Science and Engineering} {\bf 143}, Springer, Cham, 2021, 13--36.

\bibitem{FKMN1} V.O. Manturov, D. Fedoseev, S. Kim, I. Nikonov, On groups $G_n^k$ and $\Gamma_n^k$: A study of manifolds, dynamics, and invariants, \emph{Bull. Math. Sci.} {\bf 11}:2 (2021) 2150004.


\bibitem{M22} V.O. Manturov, Realisability of $G^3_n$, realisability projection, and kernel of the $G^3_n$-braid presentation, arXiv:2210.13338.



\end{thebibliography}

\end{document}

\pagebreak

The key idea is the following. We consider an $n$-component link where $n$ is an integer number greater than $3$ as a generic link in $\R^3$.
Here {\em genericity} is related to triples of points on different components:
we require that every horisontal trisecant does not contain any other point of the link, does not pass through any maximal
or minimal point of the link and is {\em non-degenerate} in the sence of \cite{MN}. 

Now, we pass to MN-indices.

For three numbers $i$, $j$ and $k$ pairwise distinct numbers from $\{1\dots n\}$, and for a number $l\in \{1,\cdots, n\}\backslash \{i,j,k\}$,
we shall define a Manturov-Nikonov indexes for triples for horisontal trisecant such that three points belong to $i$, $j$ and $k$ components respectively.
This horszontal trisecant passes through three components and by definition it does not pass through any other point on the link.

The $l$-index of an $(i,j,k)$-quadrisecant is valued in a pair of opposite vertices of the discrete cube $\{0,1\}^{3}$,
so, it may have one of the four possible values. In \cite{MN}, such indices were used for constructing lots of invariants of 
braids valued in the free product $\Z_{2}*\Z_{2}*\Z_{2}*\Z_{2}$.

Here we define the index as follows. Given a horisontal trisecant of the link $L=L_{1}\sqcup \dots \sqcup L_{n}$ and a generic horisontal
trisecant of it.

\begin{dfn}
Now the index for a trisecant $z_i$, $z_j$, $z_k$ belonging to the same horisontal line is defined as follows. We take our index $l$ and count the following three things, modulo $2$.
We take all possible points $w$, $l$ belonging to $l$ component and for on the horizontal plane containing $z_i$, $z_j$, $z_k$ and whatever $w$, $l$, we count the number of triangles modulo 2,
$z_i,z_j,w_l$, $z_i,z_k,w_l$ and $z_j,z_k,w_l$ which are positively oriented and take these three numbers mod $2$.
\end{dfn}


This agrees with our definition from \cite{MN} for {\em braids} $w_l$ and it switches well to the opposite if we pass through a quadrisecant $z_iz_jz_k$ and whatever point on $l$ component.

Consider a link $L=L_{1}\sqcup \cdots\sqcup L_{n}$ represented as a closure of a braid $\beta$. Let us realise it as follows.

$n$ points on the top plane and $n$ points on the bottom plane under them.

\begin{lm}
Those coming from local maxima and local minima coincide.
\end{lm}


\section{Constructing words from a braid.}

Let $G$ be the group $\Z_{2}*\Z_{2}*\Z_{2}*\Z_{2}$ with generators
$a_{ijk}^{000},a_{ijk}^{001},a_{ijk}^{010},a_{ijk}^{011}$ 
which we shall denote for simplicity by $a,b,c,d$.

The aim of the present section is to construct a word in these four generators.

We fix three strands corresponding to components $i,j,k$. Each strand can be either the strand of the braid $\beta$ or
a vertical closing strand going from the top to the bottom.

We start scanning our link by horisontal plane from the top (maximum level) to the bottom (minimum level) by horisontal planes.

For each horisontal plane containing an $(i,j,k)$-trisecant we write the corresponding letter with index.
Say, if our trisecant has index $(0,0,0)$ or $(1,1,1)$ we write the generator $z_{i,j,k}^{0,0}\sim a$,
for $(0,1,0)$ or $(1,0,1)$ we write $z_{i,j,k}^{0,1}\sim b$,
for $(1,0,0)$ or $(0,1,1)$ we write $z_{i,j,k}^{1,0}\sim c$,
and for

\subsection{
Definition of a collection of words}

What we have done above agrees with the construction of \cite{MN} with the following two differences:
\begin{enumerate}
\item In \cite{MN} we considered braids, where each strand goes from the top to the bottom, and 
there is no need to choose any concrete strands 
\item On each horisontal plane there are more than one point belonging to each of the components.
\end{enumerate}

What to do with this situation? 

We shall take all possible choices of triples of lines. Namely, if our braid $\beta$ contains $m_{1}$ strands belonging to 
component 1, $m_{2}$ strands belonging to component 2, and $m_{3}$ strands belonging to component 3,
then we well have $M=8 m_{1} m_{2} m_{3}$ such triples of lines 
(here 8 appears because for each strand of $\beta$ we have two strands of the knot).

Hence, we have a collection of $M$ words in $a,b,c,d$.
Denote the corresponding conjugacy classes of words by $W_{1}(\beta),\cdots, W_{M}(\beta)$.


\section{The first Markov move}

\begin{thm}
If two braids $\beta$ and $\beta'$ are conjugate (hence represent the same
$n$-component link $L$, then the corresponding words are conjugate,
hence $W_{j}(\beta)$ coincides with $W_{j}(\beta')$.
\label{th1}
\end{thm}

\section{The second Markov move}

Again, let $\beta$ be a braid representing the link $L=L_{1}\sqcup\cdots\sqcup L_{n}$.
Let $\beta'$ be obtained from $\beta$ by a stabilisation (negative or positive).

This stabilisation may concern either one of the stands $1,2,3$ or one of the 
remaining strands. In the latter case, the number of words $M$ does not change.
In the first case, say, for the first strand we have $m'_{1}= m_{1}+1$,
hence, the number $M'$ for $\beta$ is equal to $M=8 (m_{1}+1)m_{2}m_{3}$.

Hence,words for $\beta'$ contains $8m_{1}m_{2}m_{3}$ words corresponding to $\beta$
together with $4m_{1}m_{2}$ {\em pairs of new words.}

Each pair correspond to one new strand of $\beta'$ and its closing (vertical) strand
(with all strands of components $2,3$ being the same). Hence, we get some pairs
of {\em new conjugacy classes}

\begin{thm}
In the event of conjugation of a strand $4,\cdots, n$ the conjugacy classes $W_{1},W_{2},W_{3},W_{4}$ for
$\beta$ correspond to those for $\beta'$.

If we conjugate one of the strands $1,2,3$ then
the corresponding conjugacy classes for $\beta$ and $\beta'$ agree. 

New conjugacy classes for $\beta'$ belonging to the same pair coincide. 
\label{th2}
\end{thm}



\section{The main result}

Let $\Z_{2} G$ be the set of $\Z_{2}$-linear combinations of words in $\Z_{2}*\Z_{2}*\Z_{2}*\Z_{2} =\langle a,b,c,d|a^{2}=b^{2}=c^{2}=d^{2}=1\rangle$.

For the link $L$ represented as above by the closure of $\beta$, let 
$w(\beta)$ be the sum of all conjugacy classes $W$ over all triples of strands as above.




Taking together the above theorems, we get the main result of the paper.
\begin{thm}
$w(\beta)$ is an invariant of $n$-component links.
\end{thm}

\begin{proof}
Indeed, when we apply the isotopy to $\beta$, all corresponding conjugacy classes remain the same.
The proof agrees with that for braids from \cite{MN}.

When we apply the conjugacy to $\beta$, the invariance follows from Theorem \ref{th1},
and for the second Markov move the invariance follows from Theorem \ref{th2}.

\end{proof}



\section{Calculations}

\section{Further directions}

